% ME563 - HW 3
% Tyler Jones

\documentclass[11pt]{article}
\usepackage[margin=1in]{geometry}
\usepackage{amsmath, amssymb}
\usepackage{mathpazo}
\usepackage[T1]{fontenc}
\usepackage{microtype}
\usepackage{graphicx}
\usepackage{booktabs}
\usepackage{float}
\usepackage{xcolor}
\usepackage[font=small, labelfont=bf]{caption}
\usepackage{titlesec}
\usepackage{fancyhdr}
\usepackage[colorlinks=true, allcolors=blue]{hyperref}

% ---- headers / footers ------------------------------------------------
\pagestyle{fancy}
\fancyhf{}
\setlength{\headheight}{14pt}
\renewcommand{\headrulewidth}{0.4pt}
\fancyhead[L]{\small ME 563: HW 3}
\fancyhead[R]{\small Tyler Jones}
\fancyfoot[C]{\small \thepage}

% ---- section spacing --------------------------------------------------
\titleformat{\section}{\large\bfseries}{}{0pt}{}[\vspace{-0.6\baselineskip}\rule{\textwidth}{0.4pt}]
\titlespacing*{\section}{0pt}{1.6\baselineskip}{0.8\baselineskip}
\titlespacing*{\subsection}{0pt}{1.1\baselineskip}{0.4\baselineskip}

% ---- table spacing ----------------------------------------------------
\renewcommand{\arraystretch}{1.25}

% ---- shorthands -------------------------------------------------------
\newcommand{\pd}[2]{\frac{\partial #1}{\partial #2}}
\newcommand{\vect}[1]{\mathbf{#1}}

\title{\vspace{-1.5em}\textbf{ME563 HW 3}}
\author{Tyler Jones}
\date{Due: September 30, 2026}

\begin{document}
\maketitle
\thispagestyle{fancy}

%=======================================================================
\section*{Problem 1}

\noindent\textbf{Problem statement.} The components of the mass flow vector $\rho\vec u$ are
$\rho u_x = 4x^2y$, $\rho u_y = xyz$, and $\rho u_z = yz^2$.

\subsection*{(a)}

\noindent\textit{Problem statement.} Compute the net mass outflow through the closed surface formed
by the planes, $x=0$, $x=1$, $y=0$, $y=1$, $z=0$, $z=1$.

\medskip
\noindent $\rho\vect u = \langle 4x^2y,\; xyz,\; yz^2\rangle$, $\Phi = \oint_S \rho\vect u\cdot\hat{\vect n}\,dA$:
\begin{align*}
x=0,\ \hat{\vect n}=-\hat{\vect e}_x:&\quad \Phi = \int_0^1\!\!\int_0^1 -4x^2y\big|_{x=0}\,dz\,dy = 0 \\
x=1,\ \hat{\vect n}=+\hat{\vect e}_x:&\quad \Phi = \int_0^1\!\!\int_0^1 4y\,dz\,dy = 4\int_0^1 y\,dy = 2 \\
y=0,\ \hat{\vect n}=-\hat{\vect e}_y:&\quad \Phi = \int_0^1\!\!\int_0^1 -xyz\big|_{y=0}\,dz\,dx = 0 \\
y=1,\ \hat{\vect n}=+\hat{\vect e}_y:&\quad \Phi = \int_0^1\!\!\int_0^1 xz\,dz\,dx = \tfrac12\int_0^1 x\,dx = \tfrac14 \\
z=0,\ \hat{\vect n}=-\hat{\vect e}_z:&\quad \Phi = \int_0^1\!\!\int_0^1 -yz^2\big|_{z=0}\,dy\,dx = 0 \\
z=1,\ \hat{\vect n}=+\hat{\vect e}_z:&\quad \Phi = \int_0^1\!\!\int_0^1 y\,dy\,dx = \tfrac12\int_0^1 dx = \tfrac12
\end{align*}
\begin{equation*}
\boxed{\oint_S \rho\vect u\cdot\hat{\vect n}\,dA = 0 + 2 + 0 + \tfrac14 + 0 + \tfrac12 = \tfrac{11}{4}}
\end{equation*}

\subsection*{(b)}

\noindent\textit{Problem statement.} Compute $\vec\nabla\cdot(\rho\vec u)$ and integrate over the
volume bounded by the surface defined in part a).

\begin{equation*}
\nabla\cdot(\rho\vect u) = \pd{}{x}(4x^2y) + \pd{}{y}(xyz) + \pd{}{z}(yz^2)
= \boxed{8xy + xz + 2yz}
\end{equation*}
\begin{align*}
\int_0^1\!\!\int_0^1\!\!\int_0^1 (8xy + xz + 2yz)\,dz\,dy\,dx
&= \int_0^1\!\!\int_0^1 \Big(8xyz + \tfrac12 xz^2 + yz^2\Big)\Big|_{z=0}^{z=1} dy\,dx \\
&= \int_0^1\!\!\int_0^1 \Big(8xy + \tfrac12 x + y\Big)\,dy\,dx \\
&= \int_0^1 \Big(4xy^2 + \tfrac12 xy + \tfrac12 y^2\Big)\Big|_{y=0}^{y=1} dx
 = \int_0^1 \Big(4x + \tfrac12 x + \tfrac12\Big)\,dx \\
&= \Big(2x^2 + \tfrac14 x^2 + \tfrac12 x\Big)\Big|_{x=0}^{x=1}
\end{align*}
\begin{equation*}
\boxed{\int_V \nabla\cdot(\rho\vect u)\,dV = \tfrac{11}{4}}
\end{equation*}

\subsection*{(c)}

\noindent\textit{Problem statement.} Explain why the results in a) and b) should, or should not, be
equal.

\medskip
\noindent They should be equal. $\rho\vect u$ is continuously differentiable on the cube $V$, so the divergence theorem applies:
\begin{equation*}
\int_V \nabla\cdot(\rho\vect u)\,dV = \oint_S \rho\vect u\cdot\hat{\vect n}\,dA = \tfrac{11}{4}.
\end{equation*}

%=======================================================================
\section*{Problem 2}

\noindent\textbf{Problem statement.} For the following flows find the missing velocity component
needed for the flow to satisfy the incompressible continuity equation

\medskip
\noindent Incompressible continuity: $\nabla\cdot\vect u = \pd{u_1}{x_1} + \pd{u_2}{x_2} + \pd{u_3}{x_3} = 0$.
In each case $f$ is an arbitrary function of the two coordinates not differentiated against.

\subsection*{(a)}

\noindent\textit{Problem statement.} $u_1 = x_1^2 + x_2^2 + a^2, \quad v_2 = -x_1x_2 - x_2x_3 - x_1x_3, \quad u_3 = ?$

\medskip
\noindent ($v_2$ taken as $u_2$. I think this was a typo in the problem statement.)
\begin{gather*}
\pd{}{x_1}\big(x_1^2+x_2^2+a^2\big) + \pd{}{x_2}\big(-x_1x_2-x_2x_3-x_1x_3\big) + \pd{u_3}{x_3} = 0 \\
\Rightarrow\; 2x_1 - x_1 - x_3 + \pd{u_3}{x_3} = 0
\;\Rightarrow\; \pd{u_3}{x_3} = x_3 - x_1
\end{gather*}
\begin{equation*}
\int du_3 = \int (x_3 - x_1)\,dx_3
\;\Rightarrow\;
\boxed{u_3 = \tfrac12 x_3^2 - x_1x_3 + f(x_1,x_2)}
\end{equation*}

\subsection*{(b)}

\noindent\textit{Problem statement.} $u_1 = \ln\!\big(x_2^2 + x_3^2\big), \quad u_2 = \sin\!\big(x_1^2 + x_3^2\big), \quad u_3 = ?$

\begin{equation*}
\pd{}{x_1}\Big[\ln\!\big(x_2^2+x_3^2\big)\Big] + \pd{}{x_2}\Big[\sin\!\big(x_1^2+x_3^2\big)\Big] + \pd{u_3}{x_3} = 0
\;\Rightarrow\; 0 + 0 + \pd{u_3}{x_3} = 0
\;\Rightarrow\;
\boxed{u_3 = f(x_1,x_2)}
\end{equation*}

\subsection*{(c)}

\noindent\textit{Problem statement.} $u_1 = ?, \quad
u_2 = \dfrac{x_2}{\big(x_1^2+x_2^2+x_3^2\big)^{3/2}}, \quad
u_3 = \dfrac{x_3}{\big(x_1^2+x_2^2+x_3^2\big)^{3/2}}$

\medskip
\noindent Let $r = (x_1^2+x_2^2+x_3^2)^{1/2}$, so $\partial r/\partial x_i = x_i/r$ and
\begin{equation*}
\pd{}{x_i}\left(\frac{x_i}{r^3}\right) = \frac{1}{r^3} - \frac{3x_i}{r^4}\frac{x_i}{r} = \frac{r^2 - 3x_i^2}{r^5}
\quad(\text{no sum}).
\end{equation*}
Continuity:
\begin{equation*}
\pd{u_1}{x_1} + \frac{r^2 - 3x_2^2}{r^5} + \frac{r^2 - 3x_3^2}{r^5} = 0
\;\Rightarrow\;
\pd{u_1}{x_1} = \frac{3\big(x_2^2+x_3^2\big) - 2r^2}{r^5} = \frac{3\big(r^2 - x_1^2\big) - 2r^2}{r^5} = \frac{r^2 - 3x_1^2}{r^5}
\end{equation*}
which is $\partial(x_1/r^3)/\partial x_1$. Hence
\begin{equation*}
\int du_1 = \int \frac{r^2 - 3x_1^2}{r^5}\,dx_1
\;\Rightarrow\;
\boxed{u_1 = \frac{x_1}{\big(x_1^2+x_2^2+x_3^2\big)^{3/2}} + f(x_2,x_3)}
\end{equation*}

%=======================================================================
\section*{Problem 3}

\noindent\textbf{Problem statement.} Consider the velocity field $\vect u = (\alpha x/t, 0, 0)$ (in
Cartesian), where $\alpha$ is a constant. The density of the fluid is spatially uniform, but
time-dependent, i.e., $\rho$ is a function of $t$.

\subsection*{(a)}

\noindent\textit{Problem statement.} Determine the density field $\rho(t)$, given $\rho = \rho_0$ at
$t = t_0$.

\medskip
\noindent $\vect u = \begin{bmatrix} \alpha x/t & 0 & 0\end{bmatrix}^{\top}$, $\alpha\in\mathbb{R}$, $\rho = \rho(t)$.
Since $\nabla\rho = \vect 0$, continuity reduces to the form discussed in lecture
\begin{equation*}
\pd{\rho}{t} + \nabla\cdot(\rho\vect u) = 0
\;\Rightarrow\;
\pd{\rho}{t} + \rho\left[\pd{}{x}\left(\frac{\alpha x}{t}\right) + 0 + 0\right] = 0
\;\Rightarrow\;
\frac{d\rho}{dt} = -\rho\,\frac{\alpha}{t}
\end{equation*}
\begin{equation*}
\int_{\rho_0}^{\rho} \frac{d\rho'}{\rho'} = -\alpha\int_{t_0}^{t}\frac{dt'}{t'}
\;\Rightarrow\;
\ln\frac{\rho}{\rho_0} = -\alpha\ln\frac{t}{t_0} = \ln\left(\frac{t_0}{t}\right)^{\alpha}
\;\Rightarrow\;
\boxed{\rho(t) = \rho_0\left(\frac{t_0}{t}\right)^{\alpha}}
\end{equation*}

\subsection*{(b)}

\noindent\textit{Problem statement.} What are the temporal ($\partial\vect u/\partial t$), advective
($\vect u\cdot(\nabla\vect u)$), and total ($D\vect u/Dt$) accelerations for this flow field? What
does $\alpha = 1$ imply?

\medskip

\medskip
\noindent Temporal:
\begin{equation*}
\pd{\vect u}{t} = \pd{}{t}\left(\frac{\alpha x}{t}\right)\hat{\vect e}_x
\;\Rightarrow\;
\boxed{\pd{\vect u}{t} = -\frac{\alpha x}{t^2}\,\hat{\vect e}_x}
\end{equation*}
Advective:
\begin{equation*}
\vect u\cdot\nabla\vect u = u_x\pd{u_x}{x}\,\hat{\vect e}_x = \frac{\alpha x}{t}\cdot\frac{\alpha}{t}\,\hat{\vect e}_x
\;\Rightarrow\;
\boxed{\vect u\cdot\nabla\vect u = \frac{\alpha^2 x}{t^2}\,\hat{\vect e}_x}
\end{equation*}
Total:
\begin{equation*}
\frac{D\vect u}{Dt} = \pd{\vect u}{t} + \vect u\cdot\nabla\vect u
\;\Rightarrow\;
\boxed{\frac{D\vect u}{Dt} = \frac{\alpha(\alpha-1)\,x}{t^2}\,\hat{\vect e}_x}
\end{equation*}

\noindent $\alpha = 1$:
\begin{equation*}
\pd{\vect u}{t} = -\frac{x}{t^2}\hat{\vect e}_x, \qquad
\vect u\cdot\nabla\vect u = +\frac{x}{t^2}\hat{\vect e}_x, \qquad
\boxed{\frac{D\vect u}{Dt} = \vect 0}
\end{equation*}
The local and advective accelerations cancel each other out; hence, the flow field is unsteady at every point, but every
fluid particle moves at \textbf{constant velocity}.

%=======================================================================
\section*{Problem 4}

\noindent\textbf{Problem statement.} Consider a two-dimensional flow with velocity components given by
\begin{align*}
v_1 &= c\,x_1 \\
v_2 &= -c\,x_2
\end{align*}
Find expressions for the vorticity and the strain rate tensor, (i.e. calculate their components)

\medskip
\noindent Velocity gradient tensor:
\begin{equation*}
\pd{v_i}{x_j} =
\begin{bmatrix}
\partial v_1/\partial x_1 & \partial v_1/\partial x_2 \\
\partial v_2/\partial x_1 & \partial v_2/\partial x_2
\end{bmatrix}
=
\begin{bmatrix}
c & 0 \\
0 & -c
\end{bmatrix}
\end{equation*}

\noindent Vorticity, $\boldsymbol\omega = \nabla\times\vect v$, $\omega_i = \epsilon_{ijk}\,\partial_j v_k$.
For planar flow ($v_3 = 0$, $\partial/\partial x_3 = 0$), $\omega_1 = \omega_2 = 0$ and
\begin{equation*}
\omega_3 = \pd{v_2}{x_1} - \pd{v_1}{x_2} = 0 - 0
\;\Rightarrow\;
\boxed{\boldsymbol\omega = (0,\,0,\,0)}
\end{equation*}
The flow is irrotational ($R_{ij} = 0$).

\medskip
\noindent Strain rate tensor, $S_{ij} = \tfrac12\left(\pd{v_i}{x_j} + \pd{v_j}{x_i}\right)$:
\begin{align*}
S_{11} &= \pd{v_1}{x_1} = c, \qquad
S_{22} = \pd{v_2}{x_2} = -c, \qquad
S_{12} = S_{21} = \tfrac12\left(\pd{v_1}{x_2} + \pd{v_2}{x_1}\right) = 0
\end{align*}
\begin{equation*}
\boxed{S_{ij} =
\begin{bmatrix}
c & 0 \\
0 & -c
\end{bmatrix}}
\end{equation*}
$S_{ii} = c - c = 0$, so I think the flow is simplyy irrotational straining.

%=======================================================================
\section*{Problem 5}

\noindent\textbf{Problem statement.} Consider a plane Couette flow of a viscous fluid confined between
two flat plates a distance $b$ apart. At steady state the velocity distribution is $u = Uy/b$ and
$v = w = 0$, where the upper plate at $y = b$ is moving parallel to itself at speed $U$, and the
lower plate is held stationary. Find the rates of linear strain, the rate of shear strain, and
vorticity in the flow.

\medskip
\noindent The only nonzero velocity gradient is $\partial u/\partial y = U/b$:
\begin{equation*}
\pd{u_i}{x_j} =
\begin{bmatrix}
0 & U/b & 0 \\
0 & 0 & 0 \\
0 & 0 & 0
\end{bmatrix}
\end{equation*}

\noindent Rates of linear strain:
\begin{equation*}
\boxed{S_{xx} = \pd{u}{x} = 0, \qquad S_{yy} = \pd{v}{y} = 0, \qquad S_{zz} = \pd{w}{z} = 0}
\end{equation*}

\noindent Rate of shear strain:

\begin{equation*}
\boxed{\pd{u}{y} + \pd{v}{x} = 2S_{xy} = \frac{U}{b}}
\;\Rightarrow\;
S_{xy} = S_{yx} = \frac{U}{2b}, \qquad S_{xz} = S_{yz} = 0
\end{equation*}
\begin{equation*}
S_{ij} =
\begin{bmatrix}
0 & U/(2b) & 0 \\
U/(2b) & 0 & 0 \\
0 & 0 & 0
\end{bmatrix}
\end{equation*}

\noindent Vorticity:
\begin{equation*}
\boldsymbol\omega = \nabla\times\vect u
= \left(\pd{w}{y} - \pd{v}{z}\right)\hat{\vect e}_x
+ \left(\pd{u}{z} - \pd{w}{x}\right)\hat{\vect e}_y
+ \left(\pd{v}{x} - \pd{u}{y}\right)\hat{\vect e}_z
\;\Rightarrow\;
\boxed{\boldsymbol\omega = -\frac{U}{b}\,\hat{\vect e}_z}
\end{equation*}


\end{document}
